\section{Paris and Cologne: The School of Albert}

\subsection{The student itinerary}

Released to his order about 1245, Thomas was sent north for formation, and
for the next seven years his address is Albert's address. Albert of
Lauingen---``the Great'' already to his contemporaries---was the
Dominicans' most learned man, engaged in the audacious project of
digesting the whole of Aristotle and the Arabic learning for Latin
Christendom.

\witnesslimit{The exact itinerary of 1245--48 is reconstructed variously.
The 1912 Encyclopedia has Thomas go first to Cologne, arriving ``in 1244
or 1245,'' then accompany Albert to Paris in 1245; the standard modern
chronology (so the Stanford Encyclopedia) sends him to Paris 1245--48.
Both agree on what matters: from 1248 to 1252 he was at Cologne, in the
new Dominican \textit{studium generale}, as Albert's student. This study
asserts the shared core and leaves the first leg open.}

At Cologne the canonization-era Lives place the story that named him for
posterity: the big, silent southerner whom his fellow students called
``the dumb ox,'' until Albert, having heard him defend a difficult
thesis, declared---in the wording the 1912 Encyclopedia transmits---``We
call this young man a dumb ox, but his bellowing in doctrine will one day
resound throughout the world.'' No record from Cologne or Paris in
Thomas's lifetime preserves the nickname or the prophecy; the anecdote is
examined in section~\ref{sec:traditions}. What the documents do support
is its premise: Albert judged the student exceptional, and it was on
Albert's nomination, in 1252, that Thomas returned to Paris to begin the
theologian's apprenticeship.

\subsection{Bachelor of the Sentences}

From 1252 to 1256 Thomas lectured at Paris as a bachelor, first on
Scripture, then on the \work{Sentences} of Peter Lombard, the standard
theological textbook of the schools. His commentary---the
\work{Scriptum super Sententiis}, his first major work---already shows
the method that would define him: the calm division of every question,
the fair statement of objections, the appeal to Aristotle as
philosophical common coin, the refusal of rhetorical violence. In these
same years he wrote for his brethren the small treatise \work{De ente et
essentia}, on being and essence, whose metaphysical distinctions became
foundations of his system. The impersonality is total: nowhere in
thousands of pages does the young master record a feeling about the
family that had imprisoned him or the university that was about to try
to bar his way.
